Deux cadres se partagent ce chapitre. Les énoncés qui parlent de coordonnées vivent dans le
plan ou l'espace repérés, où un point \emph{est} la donnée de ses coordonnées ; les autres —
produit scalaire, Al-Kashi, théorème de la médiane — sont écrits sans repère, comme le fait le
programme lorsqu'il raisonne sur des vecteurs.

La plus jolie démonstration du chapitre est celle de l'aire d'un triangle. La bibliothèque
donne l'identité $\langle u, v\rangle^2 + \det(u,v)^2 = \|u\|^2\|v\|^2$, forme algébrique de
$\cos^2 + \sin^2 = 1$ pour deux vecteurs ; en y remplaçant le produit scalaire par
$\|u\|\|v\|\cos\theta$, il reste $|\det(u,v)| = \|u\|\|v\|\sin\theta$, d'où la formule
$\frac{1}{2}ab\sin C$.

Une ligne du programme est une définition et non un théorème : \og{} $u \perp v$ si et
seulement si $u \cdot v = 0$ \fg{}. En Lean, l'orthogonalité \emph{est} la nullité du produit
scalaire ; le chapitre démontre donc la caractérisation qui a un contenu — le rectangle est le
parallélogramme dont les diagonales sont égales.

La géométrie de l'espace est la partie la moins complète du livre. La classification des
positions relatives est une taxinomie plus qu'un théorème ; le théorème du toit et les
intersections de plans reposent sur un calcul de dimensions — l'intersection de deux plans
distincts de l'espace est une droite parce que $2 + 2 - 3 = 1$ — que le chapitre ne mène pas.
